% =========================================================
\subsubsection{Properties of Metric Spaces}
\label{sec:properties-metric-spaces}
% =========================================================

% ---- Toolkit ----
\begin{tcolorbox}[colback=gray!6, colframe=gray!40, arc=2pt,
  left=6pt, right=6pt, top=4pt, bottom=4pt,
  title={\small\textbf{Properties of Metric Spaces --- Quick Reference}},
  fonttitle=\small\bfseries]
\small
\begin{tabular}{l l l l}
\toprule
\textbf{Label} & \textbf{Statement} & \textbf{Proof method} & \textbf{Detail} \\
\midrule
D\ref{def:diameter}
  & $\operatorname{diam}(A) := \sup\{d(x,y):x,y\in A\}$
  & Definition
  & \hyperref[def:diameter]{$\downarrow$} \\
D\ref{def:bounded-set}
  & $A$ bounded $\iff$ $\operatorname{diam}(A) < \infty$
  & Definition
  & \hyperref[def:bounded-set]{$\downarrow$} \\
T\ref{thm:monotonicity-of-diameter}
  & $A \subseteq B \Rightarrow \operatorname{diam}(A) \le \operatorname{diam}(B)$
  & Supremum over larger set
  & \hyperref[thm:monotonicity-of-diameter]{$\downarrow$} \\
D\ref{def:subspace-metric}
  & $(A, d|_A)$ is a metric space for any $A \subseteq X$
  & Definition
  & \hyperref[def:subspace-metric]{$\downarrow$} \\
D\ref{def:isolated-point-ms}
  & $x$ isolated $\iff$ $\{x\}$ is open in $X$
  & Definition
  & \hyperref[def:isolated-point-ms]{$\downarrow$} \\
D\ref{def:discrete-metric-space}
  & Discrete $\iff$ every subset is open
  & Definition
  & \hyperref[def:discrete-metric-space]{$\downarrow$} \\
P\ref{prop:discrete-metric}
  & Discrete metric $d_{\mathrm{disc}}$ makes every point isolated
  & Direct from definition
  & \hyperref[prop:discrete-metric]{$\downarrow$} \\
\bottomrule
\end{tabular}
\end{tcolorbox}

% ---- Diameter ----
\begin{tcolorbox}[colback=propbox, colframe=propborder, arc=2pt,
  left=6pt, right=6pt, top=4pt, bottom=4pt,
  title={\small\textbf{Definition (Diameter)}},
  fonttitle=\small\bfseries]
\begin{definition}[Diameter]
\label{def:diameter}
Let $(X,d)$ be a metric space and let $A \subseteq X$.
The \emph{diameter} of $A$ is
\[
\operatorname{diam}(A)
:=
\sup\{\, d(x,y) : x,y \in A \,\}.
\]
When $A = \varnothing$ the convention is $\operatorname{diam}(\varnothing) := 0$.
The \emph{diameter of $X$} itself is $\operatorname{diam}(X)$.
\end{definition}
\end{tcolorbox}

\begin{remark*}[Fully quantified form]
\[
\forall (X,d)\ \text{metric space},\;
\forall A \subseteq X,\quad
\operatorname{diam}(A)
\;=\;
\sup\,\{\, r \in \mathbb{R} \mid
  \exists\, x, y \in A,\; d(x,y) = r \,\}.
\]
\end{remark*}

\begin{tikzpicture}[
  dot/.style={circle,fill,inner sep=1.6pt},
  lbl/.style={font=\small},
  sublbl/.style={font=\scriptsize,text=gray!60!black}
]
% Blob — shifted right so there's room on left for label
\fill[blue!10]
  (1.5,0)
  .. controls (2.4,-0.5) and (4.0,-0.3) .. (4.6, 0.8)
  .. controls (5.2, 1.8) and (4.8, 3.2) .. (3.8, 3.5)
  .. controls (2.8, 3.8) and (1.2, 3.2) .. (0.8, 2.2)
  .. controls (0.4, 1.2) and (0.6, 0.4) .. (1.5, 0)
  -- cycle;
\draw[blue!55!black, dashed, line width=0.8pt]
  (1.5,0)
  .. controls (2.4,-0.5) and (4.0,-0.3) .. (4.6, 0.8)
  .. controls (5.2, 1.8) and (4.8, 3.2) .. (3.8, 3.5)
  .. controls (2.8, 3.8) and (1.2, 3.2) .. (0.8, 2.2)
  .. controls (0.4, 1.2) and (0.6, 0.4) .. (1.5, 0)
  -- cycle;
\node[lbl,text=blue!60!black,font=\small\itshape] at (4.0,0.4) {$A$};

% Points — keep x4/x5 as diameter, others well away from that chord
\coordinate (p1) at (1.8, 2.4);   % x1 — upper left
\coordinate (p2) at (3.8, 2.8);   % x2 — upper right
\coordinate (p3) at (3.0, 1.0);   % x3 — lower centre
\coordinate (p4) at (1.3, 0.5);   % x4  (diameter SW endpoint)
\coordinate (p5) at (4.5, 3.0);   % x5  (diameter NE endpoint)

% Non-diameter chords
\draw[gray!45, line width=0.7pt] (p1) -- (p2);
\draw[gray!45, line width=0.7pt] (p2) -- (p3);
\draw[gray!45, line width=0.7pt] (p1) -- (p3);

% Chord labels — positioned away from diameter
\node[sublbl,anchor=south,rotate=18]  at ($(p1)!0.5!(p2)+(0,0.1)$)  {$d(x_1,x_2)$};
\node[sublbl,anchor=west,rotate=-60]  at ($(p2)!0.5!(p3)+(0.12,0)$) {$d(x_2,x_3)$};

% Diameter chord
\draw[red!70!black, line width=2.0pt] (p4) -- (p5);

% Diameter label — outside blob, below-left of chord
\draw[red!50!black,line width=0.5pt,dashed]
  ($(p4)!0.45!(p5)$) -- ($(p4)!0.45!(p5)+(-1.2,-0.6)$);
\node[lbl,text=red!70!black,anchor=north east]
  at ($(p4)!0.45!(p5)+(-1.2,-0.6)$)
  {$d(x_4,x_5) = \operatorname{diam}(A)$};

% Points and labels
\foreach \p in {p1,p2,p3,p4,p5} { \node[dot] at (\p) {}; }
\node[lbl,anchor=east]       at ($(p1)+(-0.1,0)$)    {$x_1$};
\node[lbl,anchor=south]      at ($(p2)+(0,0.1)$)     {$x_2$};
\node[lbl,anchor=north]      at ($(p3)+(0,-0.1)$)    {$x_3$};
\node[lbl,anchor=north east] at ($(p4)+(-0.08,-0.1)$){$x_4$};
\node[lbl,anchor=south west] at ($(p5)+(0.08,0.1)$)  {$x_5$};

% Footer
\node[lbl,align=center] at (2.8,-1.0)
  {$\operatorname{diam}(A) = \sup\{\,d(x,y) : x,y \in A\,\}$};
\end{tikzpicture}

\begin{remark*}[Intuition]
The diameter is the supremum of all pairwise distances in $A$ ---
the "widest" the set gets.
In $\mathbb{R}$ with the standard metric, $\operatorname{diam}([a,b]) = b - a$.
In $\mathbb{R}^2$, the diameter of a closed disc of radius $r$ is $2r$,
achieved by any pair of antipodal boundary points.
\end{remark*}

% =========================================================
% Antipodal Boundary Points (Metric Space Version)
% =========================================================

\begin{definition}[Antipodal Boundary Points]
Let $(X,d)$ be a metric space and let $A \subseteq X$ be a nonempty subset.
The \emph{diameter} of $A$ is
\[
\operatorname{diam}(A)
:=
\sup\{ d(x,y) : x,y \in A \}.
\]

Two points $p,q \in \partial A$ are called \emph{antipodal boundary points}
if
\[
d(p,q) = \operatorname{diam}(A).
\]
\end{definition}

\begin{remark*}[Common error --- ``upper bound for $A$'']
It is tempting to say that $\operatorname{diam}(A)$ is an upper bound
\emph{for $A$}, but this is a type error unless $X = \mathbb{R}$.
The supremum in $\operatorname{diam}(A) = \sup\{\, d(x,y) \mid x, y \in A \,\}$
is taken over the set of \emph{real numbers}
$\{\, d(x,y) \mid x, y \in A \,\}$, not over the set of points $A \subseteq X$.
An upper bound for $\operatorname{diam}(A)$ is a real number $M$ such that
$d(x,y) \le M$ for all $x, y \in A$ --- it bounds the \emph{distances},
not the points.
The correct phrasing is always:
``$M$ is an upper bound for $\{\, d(x,y) \mid x, y \in A \,\}$,''
or equivalently, ``$d(x,y) \le M$ for all $x, y \in A$.''
\end{remark*}

\begin{example*}[Why ``upper bound for $A$'' fails in $\mathbb{R}^2$]
Let $X = \mathbb{R}^2$ with the Euclidean metric and let
$A = \{(1,0),(0,1)\} \subseteq \mathbb{R}^2$.
The pairwise distances form the set
$\{\, d(x,y) \mid x, y \in A \,\} = \{0, \sqrt{2}\}$,
so $\operatorname{diam}(A) = \sqrt{2}$.
Now consider what ``$\sqrt{2}$ is an upper bound for $A$'' would mean:
we would need $(1,0) \le \sqrt{2}$ and $(0,1) \le \sqrt{2}$, but elements
of $A$ are points in $\mathbb{R}^2$, not real numbers.
The comparison is undefined --- $\mathbb{R}^2$ carries no total order
compatible with the metric.
What \emph{is} well-typed is: $\sqrt{2}$ is an upper bound for
$\{0,\sqrt{2}\} \subseteq \mathbb{R}$, since $0 \le \sqrt{2}$ and
$\sqrt{2} \le \sqrt{2}$.
\end{example*}

% ---- Bounded Set ----
\begin{tcolorbox}[colback=propbox, colframe=propborder, arc=2pt,
  left=6pt, right=6pt, top=4pt, bottom=4pt,
  title={\small\textbf{Definition (Bounded Set)}},
  fonttitle=\small\bfseries]
\begin{definition}[Bounded Set]
\label{def:bounded-set}
Let $(X,d)$ be a metric space and let $A \subseteq X$.
The set $A$ is called \emph{bounded} if $\operatorname{diam}(A) < \infty$,
i.e.\ if there exists $M \in \mathbb{R}$ such that $d(x,y) \le M$ for all
$x, y \in A$.
\end{definition}
\end{tcolorbox}

\begin{remark*}[Fully quantified form]
\[
A\ \text{is bounded}
\iff
\exists M \in \mathbb{R}^+,\;
\forall x, y \in A,\;
d(x,y) \le M.
\]
\end{remark*}

\begin{remark*}[Equivalent characterisation]
$A$ is bounded if and only if $A$ is contained in some open ball:
$\exists x_0 \in X,\; \exists r > 0,\; A \subseteq B_r(x_0)$.
The two conditions are equivalent by the triangle inequality.
\end{remark*}

% ---- Monotonicity of Diameter ----
\begin{tcolorbox}[colback=thmbox, colframe=thmborder, arc=2pt,
  left=6pt, right=6pt, top=4pt, bottom=4pt,
  title={\small\textbf{Theorem (Monotonicity of Diameter)}},
  fonttitle=\small\bfseries]
\begin{theorem}[\texorpdfstring{%
  \hyperref[prf:monotonicity-of-diameter]{Monotonicity of Diameter}%
}{Monotonicity of Diameter}]
\label{thm:monotonicity-of-diameter}
Let $(X,d)$ be a metric space and let $A \subseteq B \subseteq X$.
Then
\[
\operatorname{diam}(A) \le \operatorname{diam}(B).
\]
\end{theorem}
\end{tcolorbox}

\begin{remark*}[Fully quantified form]
\[
\forall (X,d)\ \text{metric space},\;
\forall A, B \subseteq X,\;
A \subseteq B
\;\Rightarrow\;
\sup\{\,d(x,y):x,y\in A\,\}
\;\le\;
\sup\{\,d(x,y):x,y\in B\,\}.
\]
\end{remark*}

\begin{remark*}[Intuition]
A subset can only be as wide as its superset.
Any pair of points in $A$ is also a pair in $B$, so the supremum
over $B$ is taken over a larger collection and cannot be smaller.
\end{remark*}



% ---- Isolated Point ----
\begin{tcolorbox}[colback=propbox, colframe=propborder, arc=2pt,
  left=6pt, right=6pt, top=4pt, bottom=4pt,
  title={\small\textbf{Definition (Isolated Point)}},
  fonttitle=\small\bfseries]
\begin{definition}[Isolated Point]
\label{def:isolated-point-ms}
Let $(X,d)$ be a metric space.
A point $x \in X$ is called \emph{isolated} if the singleton set $\{x\}$
is open in $X$.
\end{definition}
\end{tcolorbox}

\begin{remark*}[Equivalent formulations]
The following are equivalent for $x \in X$:
\begin{enumerate}[label=(\roman*)]
  \item $\{x\}$ is open in $X$.
  \item $\exists r > 0$ such that $B_r(x) = \{x\}$.
  \item $x$ is not a limit point of $X$
        (see Definition~\ref{def:limit-point}).
\end{enumerate}
The equivalence of (i) and (ii) is immediate from the definition of open set.
\end{remark*}

\begin{remark*}[Source note]
This is the same concept as Definition~\ref{def:isolated-point} in the
limit points section, stated there for subsets $E \subseteq X$.
Here the definition applies to the ambient space $X$ itself,
which is the natural context for discrete spaces.
\end{remark*}

% ---- Discrete Metric Space ----
\begin{tcolorbox}[colback=propbox, colframe=propborder, arc=2pt,
  left=6pt, right=6pt, top=4pt, bottom=4pt,
  title={\small\textbf{Definition (Discrete Metric Space)}},
  fonttitle=\small\bfseries]
\begin{definition}[Discrete Metric Space]
\label{def:discrete-metric-space}
Let $(X,d)$ be a metric space.
The space is called \emph{discrete} if every subset of $X$ is open,
or equivalently if every point of $X$ is isolated.
\end{definition}
\end{tcolorbox}

\begin{remark*}[Fully quantified form]
\[
(X,d)\ \text{discrete}
\iff
\forall x \in X,\;
\exists r > 0,\;
B_r(x) = \{x\}.
\]
\end{remark*}

\begin{tikzpicture}[
  dot/.style={circle,fill,inner sep=1.8pt},
  lbl/.style={font=\small},
  sublbl/.style={font=\scriptsize,text=gray!55!black}
]
% Scattered isolated points with isolation balls r=1/2
\def\pts{
  {0.0/0.0},
  {2.2/0.3},
  {4.1/0.1},
  {1.0/1.8},
  {3.3/1.6},
  {0.5/3.2},
  {2.8/3.0},
  {4.5/2.2}
}

% Draw isolation balls first (behind dots)
\foreach \xy in {
  0.0/0.0, 2.2/0.3, 4.1/0.1,
  1.0/1.8, 3.3/1.6,
  0.5/3.2, 2.8/3.0, 4.5/2.2
}{
  \pgfmathsetmacro{\px}{0}  % dummy
  \edef\pxy{\xy}
}

\foreach \px/\py in {
  0.0/0.0, 2.2/0.3, 4.1/0.1,
  1.0/1.8, 3.3/1.6,
  0.5/3.2, 2.8/3.0, 4.5/2.2
}{
  \draw[gray!35, dotted, line width=0.8pt] (\px,\py) circle (0.52cm);
}

% Draw dots on top
\foreach \px/\py in {
  0.0/0.0, 2.2/0.3, 4.1/0.1,
  1.0/1.8, 3.3/1.6,
  0.5/3.2, 2.8/3.0, 4.5/2.2
}{
  \node[dot,fill=blue!70!black] at (\px,\py) {};
}

% Annotate one point explicitly
\coordinate (ann) at (3.3,1.6);
\draw[gray!60,->,line width=0.6pt] (ann) -- ++(0.52,0);
\node[sublbl,anchor=west] at ($(ann)+(0.54,0)$) {$r = \tfrac{1}{2}$};
\node[sublbl,anchor=north] at ($(ann)+(0,-0.56)$) {$B_{1/2}(x) = \{x\}$};

% Border of ambient space
\draw[gray!30, line width=0.6pt] (-0.8,-0.7) rectangle (5.4,3.9);
\node[lbl,font=\small\itshape] at (-0.5,3.6) {$X$};

% Footer
\node[lbl,align=center] at (2.3,-1.15)
  {$d_{\mathrm{disc}}(x,y) = \begin{cases}0 & x=y \\ 1 & x\neq y\end{cases}$
   \quad $\Rightarrow$ \quad
   every point is isolated};
\end{tikzpicture}

\begin{remark*}[Intuition]
In a discrete space no point is a limit of a non-constant sequence from $X$.
Every sequence that converges must eventually be constant.
This makes discrete spaces the ``opposite'' of dense spaces.
\end{remark*}

% ---- Discrete Metric ----
\begin{proposition}[\texorpdfstring{%
  \hyperref[prf:discrete-metric]{Discrete Metric}%
}{Discrete Metric}]
\label{prop:discrete-metric}
Let $X$ be any non-empty set and define
\[
d_{\mathrm{disc}}(x,y)
:=
\begin{cases}
0 & \text{if } x = y, \\
1 & \text{if } x \neq y.
\end{cases}
\]
Then $(X, d_{\mathrm{disc}})$ is a metric space in which every point is
isolated and every subset of $X$ is open.
\end{proposition}

\begin{remark*}[Fully quantified form]
\[
\forall x \in X,\;
B_{1/2}(x) = \{x\},
\]
so $\{x\}$ is open for every $x$, and any $S \subseteq X$ is a union
of open singletons, hence open.
\end{remark*}

\begin{remark*}[Consequence]
The discrete metric shows that discreteness is not vacuous:
it is realised by a concrete metric on any set.
It also shows that the topology of a metric space depends on the metric,
not just the underlying set: $(\mathbb{R}, d_{\mathrm{disc}})$ and
$(\mathbb{R}, |\cdot|)$ have the same points but entirely different
open sets.
\end{remark*}